\section{Introduction}\label{sec:intro}

Battery--supercapacitor hybrid energy storage systems (HESSs) are used in electric vehicles (EVs) for two purposes: to keep fast power transients away from the battery and to recover braking energy~\cite{khaligh2010,cao2012}. A HESS controller has two coupled tasks.
\begin{itemize}
\item \emph{Fast regulation} of the DC-bus voltage under an unmeasured, rapidly varying traction load.
\item \emph{Slow energy management:} deciding which source supplies which part of the power, and how much energy the supercapacitor (SC) should hold in reserve.
\end{itemize}

For the fast task, disturbance-observer-based control (DOBC)~\cite{chen2016dobc,li2014dobc} combined with sliding-mode control (SMC)~\cite{utkin2009,shtessel2014} is attractive. The observer cancels most of the load disturbance, and the switching term only has to dominate the residual estimation error. This residual is exactly where fixed-gain designs fail. In a reproduction study of a recent DO-SMC design for HESS~\cite{ha2026gtsd}, we found two problems. First, its fixed switching gain violates its own sufficient stability condition by more than an order of magnitude under an acceleration transient. Second, it then performs no better than a PI controller with the same small-signal poles (Section~\ref{sec:results}). Adaptive and barrier-function SMC gains for HESS and DC-bus regulation have been proposed~\cite{ojvt2025barrier,bagherwal2024}. These do not exploit the information that the observer already provides about its own error.

Switching-gain adaptation is also not new. In DOB-assisted servo control the gain is raised when the state leaves a preset boundary~\cite{tian2022adaptive}. Extremum seeking has been used to adapt the SMC gain of a bidirectional storage converter~\cite{trinh2023}.

For the slow task, filter-based splitting with a fixed or rule-adapted time constant is the industrial baseline~\cite{nguyen2021adaptive,corapsiz2022,access2019variabletau}. Learning-based energy management is increasingly studied:
\begin{itemize}
\item TD3 for battery/ultracapacitor allocation~\cite{liu2023td3};
\item deep Q-learning on top of fixed-gain or adaptive fuzzy SMC in DC microgrids~\cite{charfeddine2025,jouili2026};
\item proximal policy optimisation with a predictive safety filter for battery/SC allocation in EVs~\cite{xu2025safe}.
\end{itemize}
Reinforcement learning has also been used to tune controller and disturbance-observer gains, both in power converters~\cite{jiang2023drl,muktiadji2024} and in general disturbance-observer-based control~\cite{choi2024dob,luna2026}. Across these works, the safety layer (when present) protects power and energy limits. The stability certificate of the underlying voltage loop is not part of the safety argument, even when the agent changes the gains on which that certificate depends.

This paper couples the two tasks in a two-tier architecture (Fig.~\ref{fig:arch}). The safety argument is organised so that the learning agent \emph{cannot remove} the stability certificate of the fast loop. The contributions are as follows.
\begin{enumerate}
\item \textbf{A certificate-enforcing adaptive DO-SMC (A-DO-SMC).} We use the identity $\dot{\hat D}=\lambda\tilde D$ of the nonlinear disturbance observer to estimate the observer error online. The switching gain is adapted so that the reaching condition holds for any observer gain. A boundary layer that scales with the gain keeps the small-signal poles fixed, and a conditional integrator removes windup.
\item \textbf{A safety-filtered TD3 supervisor} that adapts the filter time constant, the SC energy target and the observer gain from controller-internal signals only. The filter enforces the slow-variation and bound assumptions required by tier 1, together with a power-capability reserve derived from vehicle ratings and training-set statistics.
\item \textbf{A reproducible benchmark with a failure analysis.} Standard drive cycles are split strictly into train, validation and test sets. The benchmark uses three training seeds, paired noise seeds, a non-learning SOC-adaptive baseline, filter ablations and three additional held-out cycles. It shows where the learned supervisor helps and where it fails. The failures are traced to a specific mechanism: the agent uses a set-point as a fast actuator, a filter that bounds the set-point cannot detect this, and the result is power deficits on unseen cycles.
\end{enumerate}

Learning the observer gain is not new in itself~\cite{choi2024dob,luna2026}. What distinguishes this work is how responsibility is split. The reaching condition of the voltage loop is enforced online from the observer's own error estimate, whatever gain the agent selects. The safety filter then only has to keep the operating assumptions of that certificate valid. Within a web search of 16 queries (November 2026; landing pages and abstracts only), we found no HESS study that combines these elements. A systematic database search is recommended before submission. All results in this paper are from offline simulation with an averaged converter model. The limits of this evidence are discussed in Section~\ref{sec:discussion}.

The remainder of the paper is organised in five sections:
\begin{itemize}
\item Section~\ref{sec:model} presents the HESS model, the control structure and the problem formulation.
\item Section~\ref{sec:method} develops the proposed two-tier strategy: the certificate-enforcing adaptive DO-SMC, the safety-filtered TD3 supervisor, and the complete discrete-time implementation.
\item Section~\ref{sec:results} describes the simulation setup, reports the results on the reproduced profiles and on held-out drive cycles, and discusses mechanisms and limitations.
\item Section~\ref{sec:conclusion} concludes the paper.
\end{itemize}
